% Shared preamble of the two documents (Director's Research Guide, Complete Technical Research Report).
\usepackage[margin=2.3cm]{geometry}
\usepackage[T1]{fontenc}
\usepackage{lmodern}
\usepackage[utf8]{inputenc}
\usepackage{microtype}
\usepackage{amsmath,amssymb}
\usepackage{graphicx}
\usepackage{booktabs,longtable,array,tabularx,multirow}
\usepackage{xcolor}
\usepackage{float}
\usepackage[labelfont=bf,font=small]{caption}
\usepackage{enumitem}
\usepackage{url}
\usepackage{fancyhdr}
\usepackage{listings}
\usepackage[most]{tcolorbox}
\usepackage{tikz}
\usetikzlibrary{positioning,arrows.meta,fit,shapes.geometric,calc,backgrounds}
\usepackage[colorlinks=true,linkcolor=blue!50!black,urlcolor=blue!50!black,citecolor=blue!50!black]{hyperref}

\definecolor{proposedblue}{HTML}{2A78D6}
\definecolor{baselineorange}{HTML}{EB6834}
\definecolor{ruleblack}{HTML}{8A8984}
\definecolor{softgreen}{HTML}{1BAF7A}
\definecolor{softred}{HTML}{E34948}
\definecolor{softviolet}{HTML}{6250D6}

\DeclareRobustCommand{\code}[1]{\texttt{\detokenize{#1}}}   % identifiers and file names (underscores allowed; robust for section titles)
\newcommand{\pth}[1]{\path{#1}}                              % file system paths (line-breakable)
\newcommand{\codename}[1]{}                                  % internal nicknames are no longer printed in running text

% Boxes used throughout
\newtcolorbox{plainwords}[1][In plain words]{enhanced,breakable,colback=blue!4,colframe=blue!45!black,fonttitle=\bfseries\small,title=#1,
  left=6pt,right=6pt,top=4pt,bottom=4pt,boxrule=0.6pt,arc=2pt}
\newtcolorbox{implicationbox}[1][Research implication]{enhanced,breakable,colback=green!5!white,colframe=green!35!black,fonttitle=\bfseries\small,title=#1,
  left=6pt,right=6pt,top=4pt,bottom=4pt,boxrule=0.6pt,arc=2pt}
\newtcolorbox{cautionbox}[1][Important to know]{enhanced,breakable,colback=red!4,colframe=red!50!black,fonttitle=\bfseries\small,title=#1,
  left=6pt,right=6pt,top=4pt,bottom=4pt,boxrule=0.6pt,arc=2pt}

% A figure with a descriptive caption followed by a plain-English reading and the research implication.
% arguments: 1 file name (without extension), 2 width as a fraction of the text width, 3 title, 4 what is shown, 5 plain-English reading, 6 research implication, 7 label
\newcommand{\explainedfigure}[7]{%
\begin{figure}[H]\centering
\includegraphics[width=#2\textwidth,height=0.56\textheight,keepaspectratio]{figures/#1.pdf}
\caption{\textbf{#3.} #4}\label{#7}
\end{figure}
\begin{tcolorbox}[enhanced,breakable,colback=gray!6,colframe=gray!55,boxrule=0.5pt,arc=2pt,left=6pt,right=6pt,top=3pt,bottom=3pt]
{\small\textbf{Plain-English reading.} #5\par\smallskip
\textbf{Research implication.} #6}
\end{tcolorbox}\medskip}

\lstset{basicstyle=\ttfamily\footnotesize,breaklines=true,frame=single,columns=fullflexible,keepspaces=true,
        showstringspaces=false,backgroundcolor=\color{black!3}}
\setlist{itemsep=2pt,topsep=3pt}
\setlength{\parindent}{0pt}
\setlength{\parskip}{5pt}
\sloppy
\emergencystretch=2em
\setcounter{tocdepth}{2}
\renewcommand{\arraystretch}{1.12}

% Flow diagram of one reinforcement-learning attempt: starting model -> initialisation -> training set-up -> training loop -> result.
% arguments: 1 starting model, 2 initialisation, 3 training set-up (environment and reward), 4 training loop (algorithm, steps), 5 result (what happened), 6 colour of the result box
\newcommand{\rlflow}[6]{%
\begin{center}
\begin{tikzpicture}[font=\scriptsize,>=Stealth,node distance=0.35cm,
  stage/.style={draw,rounded corners=2pt,fill=white,text width=2.75cm,minimum height=3.7cm,align=left,inner sep=3pt},
  result/.style={draw,rounded corners=2pt,text width=2.75cm,minimum height=3.7cm,align=left,inner sep=3pt,line width=0.8pt}]
\node[stage,fill=blue!6] (a) {\textbf{1. Starting model}\\[2pt] #1};
\node[stage,right=of a] (b) {\textbf{2. Initialisation}\\[2pt] #2};
\node[stage,right=of b] (c) {\textbf{3. Training set-up}\\[2pt] #3};
\node[stage,right=of c] (d) {\textbf{4. Training loop}\\[2pt] #4};
\node[result,right=of d,fill=#6!12,draw=#6!70!black] (e) {\textbf{5. What happened}\\[2pt] #5};
\foreach \x/\y in {a/b,b/c,c/d,d/e} \draw[->] (\x) -- (\y);
\end{tikzpicture}
\end{center}}
